\section{FREE LIE ALGEBRAS.}

It remains for us to speak of a series of works on Lie algebras where the link with the theory of Lie groups is very tenuous; these researches have on the other hand important applications in the theory of ``abstract'' groups and more especially nilpotent groups.

The origin of this is the work of P. Hall {[}144{]}, which appeared in 1932. There is however there no question of Lie algebras: P. Hall has in view the study of a certain class of p-groups, those that he calls ``regular''. But that leads him to examine in detail iterated commutators and the lower central series of a group; he establishes on this occasion a variant of the Jacobi identity, as well as the ``Hall formula''

\[
(x y) ^ {n} = x ^ {n} y ^ {n} (x, y) ^ {n (n - 1) / 2} \dots
\]

Soon after (in 1935-1937) the fundamental works of W. Magnus {[}215 a and b{]} and E. Witt {[}337 b{]} appear. In {[}215 a{]} Magnus uses the same algebra of formal series \(\hat{A}\) as Hausdorff (since called the ``Magnus algebra''); he embeds the free group \(F\) in it and uses the natural filtration of \(\hat{A}\) in order to obtain a decreasing sequence \((F_n)\) of subgroups of \(F\) ; it is one of the first examples of filtration. He conjectures that the \(F_n\) coincide with the terms of the lower central series of \(F\) . This conjecture is proved in his second memoir {[}215 b{]}; it is also there that he explicitly makes the link between his ideas and those of P. Hall, and that he defines the free Lie algebra \(L\) (as a sub algebra of \(A\) ) of which he shows in substance that it is identified with the graded \(F\) . In {[}337 b{]}, Witt completes this result on several points. He shows notably that the enveloping algebra of \(L\) is a free associative algebra and deduces from that immediately the rank of the homogeneous components of \(L\) (``Witt formulae'').

As for the determination of the basis of L known by the name of ``Hall basis'', it seems that it only appears for the first time in 1950, in a note of M. Hall {[}143{]}, although it is implicit in the works of P. Hall and W. Magnus quoted above.

